% -*- coding: utf-8 -*-
\chapter{总结与未来工作展望}
\label{chap:conclusion}

\section{全文工作总结}
本学位论文立足于 6G 移动通信系统物理层关键技术的重大演进需求，针对传统香农分立架构在短码长下的率失真崩溃及数字通信固有的信道悬崖效应，以北京大学 085402 通信工程硕士研究生培养规范为标杆，开展了系统性的面向高速时变信道的轻量化自适应深度信源信道联合语义通信（Deep JSCC）研究。本文的主要学术贡献与技术成果总结如下：
\begin{enumerate}
  \item \textbf{攻克了快时变信道下的信噪比失配退化死结}：
  创新性地提出了基于瞬时 CSI 门控反馈的轻量级通道门控剪枝单元（DCGU）。该机制在信噪比突变时能够以微秒级时间快速重塑传输流形维度，在保证高保真语义传递的同时彻底消除了传统神经网络因信噪比失配导致的泛化崩溃；
  \item \textbf{化解了连续语义映射与商用数字基带芯片的失配矛盾}：
  构建了端到端可微分非均匀星座矢量量化收发机结构，打通了 Deep JSCC 直接在现有商用 QAM 射频硬件与 SDR 原型机上部署的工程路径；
  \item \textbf{完成了 3GPP 标准信道下的仿真与验证闭环}：
  在 5G NR TDL-D 高多普勒时变多径信道仿真平台中，系统评估了本文方法在低信噪比与深度多径衰落恶劣工况下的重构保真度与平滑抗衰落特性，为后续硬件测试提供了可复现的基准。
\end{enumerate}

\section{未来研究工作展望}
本论文所提方案虽在理论推导与基准仿真中展现出显著优越性，但面对 6G 全场景智能互联的长远愿景，后续仍有以下三个极具学术价值的探索维度：
\begin{enumerate}
  \item \textbf{非平稳信道下的理论语义容量上界严格证明}：
  当前收敛性推导主要基于经验平滑性假设。未来可进一步借助随机微分方程（SDE）与统计物理流形理论，建立连续非线性特征投影在快时变多径信道下的紧致信道容量解析上界；
  \item \textbf{通感算一体化（ISAC）场景下的多目标联合流形设计}：
  将雷达探测感知信号与语义通信特征进行空时频维度的物理级正交共享，在同一射频波形中协同实现毫米级目标定位跟踪与低时延语义信息交互；
  \item \textbf{专用神经计算芯片 (ASIC/FPGA) 的微系统流片级实现}：
  将本文设计的 DCGU 门控逻辑与卷积分析算子固化至定制硬件 IP 核中，在实际车规级芯片平台上开展流片验证与现场外场试验。
\end{enumerate}
